\section{Thảo luận}

\subsection{Khó khăn gặp phải}
\begin{itemize}[leftmargin=2em]
    \item Dữ liệu thu thập có nhiều trường thiếu hoặc không đồng nhất giữa các nhà cung cấp.
    \item Phân phối doanh số lệch mạnh, buộc phải xử lý log để quan sát xu hướng rõ ràng hơn.
    \item Một số biến mô tả (tác giả, NXB) có nhiều biến thể tên, cần chuẩn hóa thủ công.
\end{itemize}

\subsection{Hạn chế của thiết kế hiện tại}
\begin{itemize}[leftmargin=2em]
    \item Dashboard chủ yếu mô tả, chưa đi sâu vào kiểm định nguyên nhân.
    \item Chưa có dữ liệu thời gian đủ chi tiết để phân tích xu hướng dài hạn theo mùa vụ.
    \item Tương tác dừng ở mức lọc và hover, chưa có drill-down sâu tới từng sản phẩm.
\end{itemize}

\subsection{Hướng phát triển}
\begin{itemize}[leftmargin=2em]
    \item Bổ sung dữ liệu theo thời gian và lịch sử giá để phân tích xu hướng theo mùa vụ.
    \item Thêm phân khúc người dùng hoặc nhóm khách hàng nếu có dữ liệu hành vi.
    \item Mở rộng mô hình dự báo doanh số và kiểm thử bằng dữ liệu mới.
    \item Thu thập thêm thông tin khuyến mãi theo chiến dịch để đánh giá hiệu quả.
\end{itemize}

\subsection{Kết luận ngắn}
Dashboard đã giúp nhóm trả lời các câu hỏi trọng tâm về phân phối doanh số, tác động của chiết khấu và mối liên hệ giữa rating-review với doanh số. Các kết quả này hỗ trợ quyết định về danh mục và chiến lược bán hàng, đồng thời mở ra hướng phân tích sâu hơn trong các giai đoạn tiếp theo.
